\section{Discussion}
\label{sec:discussion}

Our results show that SSIM-based constraints consistently outperform L$\infty$ constraints on MNIST. This observation can be explained by three factors. First, MNIST images are nearly binary (black background, white digits), so optimal perturbations tend to flip pixels between 0 and 1, making L$\infty$ naturally close to 1.0 regardless of the constraint strength. Second, SSIM captures structural relationships between neighboring pixels, which is more meaningful for preserving visual quality on digit images than per-pixel magnitude bounds. Third, L$\infty$ constraints conflict with the $\ell_0$ objective: bounding per-pixel magnitude forces the attack to modify more pixels to achieve the same adversarial effect, thereby increasing L0. This is evident in Table~\ref{tab:mnist_results}, where the L$\infty$ + SSIM variant has L0 = 28, significantly higher than SSIM-only (L0 = 22). We expect L$\infty$ constraints to be more effective on natural images (e.g., CIFAR-10, ImageNet) where pixel values are continuous and small per-pixel changes can be truly imperceptible.

\paragraph{Limitations.} SSIM computation adds approximately 15\% overhead per iteration compared to the original $\sigma$-zero, and multi-scale SSIM adds further overhead due to multiple down-sampling operations. The optimal weights depend on the dataset and model; our grid search provides guidelines but may not generalize to all settings, so practitioners should perform dataset-specific tuning. Finally, our experiments focus on MNIST, and extension to color images (CIFAR-10, ImageNet) requires additional validation, particularly for the effectiveness of L$\infty$ constraints on natural images. We also note that our attack assumes white-box access to model gradients; extending to black-box settings would require query-based approaches~\cite{guo2019simple, brendel2019decision} or transfer-based methods.

\paragraph{Future Work.} A natural extension is to replace SSIM with Learned Perceptual Image Patch Similarity (LPIPS)~\cite{zhang2018unreasonable} for better alignment with human perception, especially on natural images. Evaluating on CIFAR-10 and ImageNet would validate L$\infty$ effectiveness on natural images and assess the scalability of our approach to higher-resolution inputs. Recent advances in adversarial training with generated data~\cite{gowal2021improving} and improved training baselines~\cite{tramer2020free} suggest that perceptually-aware attacks could also enhance robust training pipelines. Furthermore, adversarially robust models have been shown to transfer better to downstream tasks~\cite{salman2020transfer}, motivating the integration of our perceptual constraints into robust training frameworks. Extending to targeted attack scenarios, where the adversary specifies the desired misclassification, may require different trade-offs between sparsity and perceptual quality. Finally, investigating whether perceptual constraints improve transferability to physical-world settings, where perturbations must survive environmental transformations, remains an open direction.

\section{Conclusion}
\label{sec:conclusion}

We presented Perceptually-Aware $\sigma$-zero, an enhanced sparse adversarial attack that incorporates visual perception constraints into the $\sigma$-zero optimization framework. By adding SSIM loss, L$\infty$ penalty, and TV regularization---combined with an adaptive weighting scheme---our method significantly improves the imperceptibility of adversarial examples while maintaining 100\% attack success rate and achieving even sparser perturbations than the original algorithm. Our comprehensive experiments on MNIST demonstrate that SSIM improves by up to 9.6\% (0.794 $\to$ 0.870), median $\ell_0$-norm decreases by 22\% (23 $\to$ 18), and all variants maintain 100\% attack success rate. The proposed method requires only 3--5 additional lines of code in the core loss computation, making it a simple yet powerful enhancement to the $\sigma$-zero algorithm. We believe this work opens new directions for designing perceptually-aware adversarial attacks that balance sparsity, effectiveness, and imperceptibility.
